\section{A hierarchical scaffold read by a fixed number of tokens}\label{sec:scaffold}

This section builds one fixed Wang set, the scaffold, and shows that a fixed number of tokens can read its hierarchy at every depth. Hierarchical tilings go back to Berger \cite{Berger66} and Robinson \cite{Robinson71}. Mozes showed that substitution tilings with unique derivation are enforced by local rules \cite{Mozes89}; unique derivation (recognizability) of aperiodic substitutions is due to Moss\'e \cite{Mosse92} and Solomyak \cite{Solomyak98}. We use Ollinger's $104$-tile substitution \cite{Ollinger08}, in the transcription of \cite{LeanWang}, and prove the facts we need from its tables, which are printed in \cref{app:tables}.

The plan is as follows. \Cref{scf:sub-base} recalls the $104$ tiles and proves that every tiling has one parent tiling. \Cref{scf:sub-ruler} adds two bits to every tile, one constant on columns and one constant on rows. The result is a Wang set $\mathbb A$ of $276$ letters with a substitution $\sigma$ of expansion $4$ (\cref{prop:scaffold}). The bits record the depth of each column and row modulo $2$, like a ruler. \Cref{scf:sub-ports} reads the full letter of a depth-$n$ block from $\ell=\binom{276}{2}+1$ fixed sites. \Cref{scf:sub-parser} defines anchors and builds the type, edge and point certificates, graphs of bounded tree width that force the readout (\cref{thm:parser}). \Cref{scf:sub-transfer} carries their conclusions to every existential pebble family (\cref{lem:transfer}) and shows that every such family contains anchors of every depth (\cref{prop:anchors}).

Throughout, $\mathbf e_1=(1,0)$ and $\mathbf e_2=(0,1)$. Intervals are sets of integers: $[a,b]=\{a,\dots,b\}$ and $[0,m)=\{0,\dots,m-1\}$. Wang sets, legal patterns and tilings are as in \cref{set:def-wang}. Tiles $t,t'$ \emph{match in direction $\mathbf e_1$} if $t_{\mathrm E}=t'_{\mathrm W}$, and \emph{in direction $\mathbf e_2$} if $t_{\mathrm N}=t'_{\mathrm S}$. A pattern on a finite set is \emph{extendable} if it is the restriction of a tiling.

\subsection{Ollinger's substitution}\label{scf:sub-base}

A tile of the base set is a triple $(\theta,\kappa,\omega)$ of a \emph{thin} type $\theta\in\{\mathsf a,\mathsf b,\mathsf c,\mathsf d\}$, a \emph{thick} type $\kappa$ and a \emph{black} type $\omega\in\{\mathsf A,\dots,\mathsf E\}$. A thick type is one of four corners $\angle\mathsf a,\angle\mathsf b,\angle\mathsf c,\angle\mathsf d$ or one of sixteen lines $\langle\alpha,\gamma\rangle$, with a vertical atom $\alpha\in\{\mathsf r_0,\mathsf r_2,\mathsf g_0,\mathsf g_2\}$ and a horizontal atom $\gamma\in\{\mathsf r_1,\mathsf r_3,\mathsf g_1,\mathsf g_3\}$. Each component has an endpoint $\partial_X$ on each side $X\in\{\mathrm N,\mathrm E,\mathrm S,\mathrm W\}$, given in \cref{scf:tab-edges}.

\begin{definition}[The base set]\label{scf:def-base}
$\mathbb O$ is the Wang set of the $104$ triples listed in \cref{scf:tab-alphabet}. The colour of $(\theta,\kappa,\omega)$ on side $X$ is the triple $(\partial_X\theta,\partial_X\kappa,\partial_X\omega)$.
\end{definition}

This is the $104$-tile set of the Lean development \cite{LeanWang}, a transcription of Ollinger's set \cite{Ollinger08}, in a compact presentation. No agreement with the published figures of \cite{Ollinger08} is claimed or used: the proofs below use only the tables of \cref{app:tables}. The four exclusions in \cref{scf:tab-alphabet} replace a list of $104$ cases. Distinct triples give distinct tiles. Indeed, the four thin types have distinct north endpoints. A line $\langle\alpha,\gamma\rangle$ has equal north and south endpoints $\alpha$ and east endpoint $\gamma$, while each corner has distinct north and south endpoints and the four corners have distinct north endpoints. The five black types have distinct bit tuples.

\begin{lemma}[Thin phase]\label{scf:thin}
Let $f$ be a tiling of $\mathbb O$. There is a unique coset $c+2\Z^2$ such that the thin type of $f(c+2v+\varsigma)$ is $\mathsf a,\mathsf c,\mathsf d,\mathsf b$ for $\varsigma=(0,0),(1,0),(0,1),(1,1)$ respectively, for every $v\in\Z^2$.
\end{lemma}

\begin{proof}
The east endpoints of the thin types are distinct, and so are the west endpoints. By \cref{scf:tab-edges}, the thin type of the east neighbour is therefore determined: $\mathsf a\leftrightarrow\mathsf c$ and $\mathsf b\leftrightarrow\mathsf d$. In the same way the north neighbour is determined: $\mathsf a\leftrightarrow\mathsf d$ and $\mathsf b\leftrightarrow\mathsf c$. These two involutions commute; their composite exchanges $\mathsf a\leftrightarrow\mathsf b$ and $\mathsf c\leftrightarrow\mathsf d$. Hence the thin type at $u+(x_1,x_2)$ is obtained from the thin type at $u$ by applying the first involution $x_1$ times and the second $x_2$ times. The sites of thin type $\mathsf a$ form one coset $c+2\Z^2$, and the other three types sit as stated.
\end{proof}

\begin{definition}[The substitution]\label{scf:def-phi}
For a triple $t=(\theta,\kappa,\omega)$, the pattern $\varphi(t)$ on $[0,2)^2$ is given by \cref{scf:tab-subst} (left). It uses the maps $F$, $\operatorname{se}$ and $\operatorname{nw}$ of \cref{app:tables}. For a pattern $w$ on a set $U\subseteq\Z^2$, the pattern $\varphi(w)$ on $2U+[0,2)^2$ is $\varphi(w)(2v+\varsigma)=\varphi(w(v))(\varsigma)$.
\end{definition}

The definition makes sense for every triple of a thin type, a thick type and a black type, whether or not it lies in $\mathbb O$. We call such a triple a \emph{formal triple}.

\begin{lemma}[The substitution is sound]\label{scf:base}
\begin{itemize}
\item[(a)] For every $t\in\mathbb O$, $\varphi(t)$ is a legal pattern of tiles of $\mathbb O$.
\item[(b)] For every pattern $w$ of tiles of $\mathbb O$ on a set $U\subseteq\Z^2$, $\varphi(w)$ is legal if and only if $w$ is legal.
\item[(c)] $\varphi$ is injective on formal triples.
\end{itemize}
\end{lemma}

\begin{proof}
(a) Let $t=(\theta,\kappa,\omega)\in\mathbb O$. Child $(0,0)$ is $(\mathsf a,\kappa,\omega)$. It lies in $\mathbb O$ because every pair $(\kappa,\omega)$ allowed for some thin type is allowed for $\mathsf a$ (\cref{scf:tab-alphabet}). Child $(1,1)$ is $(\mathsf b,F(\theta),\mathsf A)$, a corner with black $\mathsf A$, which is allowed. Child $(1,0)$ is the line $\langle\partial_{\mathrm S}\theta,\partial_{\mathrm E}\kappa\rangle$ with thin $\mathsf c$ and black $\operatorname{se}(\omega)\in\{\mathsf B,\mathsf C\}$. It is excluded only if $\operatorname{se}(\omega)=\mathsf B$ and $\partial_{\mathrm S}\theta=\mathsf r_0$, or $\operatorname{se}(\omega)=\mathsf C$ and $\partial_{\mathrm S}\theta=\mathsf r_2$. Now $\partial_{\mathrm S}\theta=\mathsf r_0$ only for $\theta=\mathsf d$, where $\omega\in\{\mathsf D,\mathsf E\}$ and $\operatorname{se}(\omega)=\mathsf C$; and $\partial_{\mathrm S}\theta=\mathsf r_2$ only for $\theta=\mathsf b$, where $\omega=\mathsf A$ and $\operatorname{se}(\omega)=\mathsf B$. Child $(0,1)$ is the line $\langle\partial_{\mathrm N}\kappa,\partial_{\mathrm W}\theta\rangle$ with thin $\mathsf d$ and black $\operatorname{nw}(\omega)\in\{\mathsf D,\mathsf E\}$. It is excluded only if $\operatorname{nw}(\omega)=\mathsf D$ and $\partial_{\mathrm W}\theta=\mathsf r_3$, which needs $\theta=\mathsf b$ and $\omega=\mathsf A$, where $\operatorname{nw}(\omega)=\mathsf E$; or $\operatorname{nw}(\omega)=\mathsf E$ and $\partial_{\mathrm W}\theta=\mathsf r_1$, which needs $\theta=\mathsf c$ and $\omega\in\{\mathsf B,\mathsf C\}$, where $\operatorname{nw}(\omega)=\mathsf D$. So all four children lie in $\mathbb O$.

The four internal edges match. The thin pattern $\mathsf a,\mathsf c,\mathsf d,\mathsf b$ matches by the thin matching rules listed in \cref{app:tables}. For the thick components, the edge from $(0,0)$ to $(1,0)$ carries $\partial_{\mathrm E}\kappa$ on both sides, the edge from $(0,0)$ to $(0,1)$ carries $\partial_{\mathrm N}\kappa$, the edge from $(1,0)$ to $(1,1)$ carries $\partial_{\mathrm S}\theta=\partial_{\mathrm S}F(\theta)$, and the edge from $(0,1)$ to $(1,1)$ carries $\partial_{\mathrm W}\theta=\partial_{\mathrm W}F(\theta)$. Here we use that $F(\theta)$ has the four endpoints of $\theta$. For the black components, $\partial_{\mathrm W}\operatorname{se}(\omega)=\partial_{\mathrm E}\omega$, $\partial_{\mathrm S}\operatorname{nw}(\omega)=\partial_{\mathrm N}\omega$, and the two edges at child $(1,1)$ carry $\partial_{\mathrm N}\operatorname{se}(\omega)=1=\partial_{\mathrm S}\mathsf A$ and $\partial_{\mathrm E}\operatorname{nw}(\omega)=1=\partial_{\mathrm W}\mathsf A$.

(b) By (a) the edges inside each block $\varphi(w(v))$ match. Let $v,v+\mathbf e_1\in U$ with $w(v)=(\theta,\kappa,\omega)$ and $w(v+\mathbf e_1)=(\theta',\kappa',\omega')$. Two fine edges cross the seam between their blocks. The lower one joins child $(1,0)$ of the first block to child $(0,0)$ of the second. Its thin endpoints are $\mathsf g_3$ and $\mathsf g_3$; its thick endpoints are $\partial_{\mathrm E}\kappa$ and $\partial_{\mathrm W}\kappa'$; its black bits are $\partial_{\mathrm E}\operatorname{se}(\omega)=\partial_{\mathrm E}\omega$ and $\partial_{\mathrm W}\omega'$. The upper one joins child $(1,1)$ to child $(0,1)$. Its thin endpoints are $\mathsf g_1$ and $\mathsf g_1$; its thick endpoints are $\partial_{\mathrm E}F(\theta)=\partial_{\mathrm E}\theta$ and $\partial_{\mathrm W}\theta'$; its black bits are $\partial_{\mathrm E}\mathsf A=0$ and $\partial_{\mathrm W}\operatorname{nw}(\omega')=0$. So both fine edges match if and only if the three components of the coarse edge match. For $v,v+\mathbf e_2\in U$ the argument is the same: the left fine edge joins child $(0,1)$ to child $(0,0)$ and records $\partial_{\mathrm N}\kappa=\partial_{\mathrm S}\kappa'$ and $\partial_{\mathrm N}\operatorname{nw}(\omega)=\partial_{\mathrm N}\omega=\partial_{\mathrm S}\omega'$, with thin endpoints $\mathsf g_2,\mathsf g_2$; the right fine edge joins child $(1,1)$ to child $(1,0)$ and records $\partial_{\mathrm N}\theta=\partial_{\mathrm S}\theta'$, with thin endpoints $\mathsf g_0,\mathsf g_0$ and black bits $0,0$.

(c) Child $(0,0)$ shows $\kappa$ and $\omega$, and child $(1,1)$ shows $F(\theta)$. The map $F$ is a bijection from thin types to corners.
\end{proof}

The proof of (b) compares components, so no list of the $104^2$ pairs is needed.

\begin{lemma}[Parents]\label{scf:parents}
Let $w$ be a legal pattern of tiles of $\mathbb O$ on $[0,2)^2$ whose thin types are $\mathsf a,\mathsf c,\mathsf d,\mathsf b$ at $(0,0),(1,0),(0,1),(1,1)$. Then $w=\varphi(t)$ for exactly one formal triple $t$. Exactly $128$ formal triples arise in this way: the $104$ tiles of $\mathbb O$ and $24$ others (\cref{scf:tab-parents}). If $w$ extends to a legal pattern on $[0,2)^2\cup\{(0,-1),(1,-1),(-1,0),(-1,1)\}$, then $t\in\mathbb O$.
\end{lemma}

\begin{proof}
The tile $w(1,1)$ has thin type $\mathsf b$, so it is a corner with black $\mathsf A$; let $\theta=F^{-1}$ of that corner. Let $\kappa$ and $\omega$ be the thick and black types of $w(0,0)$, and $t=(\theta,\kappa,\omega)$. The tile $w(1,0)$ has thin type $\mathsf c$, so it is a line with black $\mathsf B$ or $\mathsf C$. Its west and north thick endpoints must equal $\partial_{\mathrm E}\kappa$ and $\partial_{\mathrm S}F(\theta)=\partial_{\mathrm S}\theta$; so it is the line $\langle\partial_{\mathrm S}\theta,\partial_{\mathrm E}\kappa\rangle$. Its black type has west bit $\partial_{\mathrm E}\omega$, and the west bits of $\mathsf B$ and $\mathsf C$ are $0$ and $1$; so it is $\operatorname{se}(\omega)$. In the same way $w(0,1)$ is $(\mathsf d,\langle\partial_{\mathrm N}\kappa,\partial_{\mathrm W}\theta\rangle,\operatorname{nw}(\omega))$, because the south bits of $\mathsf D$ and $\mathsf E$ are $1$ and $0$. Hence $w=\varphi(t)$, and $t$ is unique by \cref{scf:base}(c).

Which $t$ occur? Child $(0,0)$ must lie in $\mathbb O$, which leaves exactly the $52$ pairs $(\kappa,\omega)$ of thin type $\mathsf a$ in \cref{scf:tab-alphabet}. Child $(1,1)$ is always allowed. Child $(1,0)$ is allowed unless $\operatorname{se}(\omega)=\mathsf B$ and $\partial_{\mathrm S}\theta=\mathsf r_0$, or $\operatorname{se}(\omega)=\mathsf C$ and $\partial_{\mathrm S}\theta=\mathsf r_2$. For $\theta=\mathsf b$ ($\partial_{\mathrm S}\theta=\mathsf r_2$) this excludes $\omega\in\{\mathsf C,\mathsf D,\mathsf E\}$; for $\theta=\mathsf d$ ($\partial_{\mathrm S}\theta=\mathsf r_0$) it excludes $\omega\in\{\mathsf A,\mathsf B\}$. Child $(0,1)$ is allowed unless $\operatorname{nw}(\omega)=\mathsf D$ and $\partial_{\mathrm W}\theta=\mathsf r_3$, or $\operatorname{nw}(\omega)=\mathsf E$ and $\partial_{\mathrm W}\theta=\mathsf r_1$. For $\theta=\mathsf b$ this excludes $\omega\in\{\mathsf B,\mathsf C,\mathsf D\}$; for $\theta=\mathsf c$ it excludes $\omega\in\{\mathsf A,\mathsf E\}$. For $\theta=\mathsf a$ the new atoms $\mathsf g_2$ and $\mathsf g_3$ are green and exclude nothing. This gives the second column of \cref{scf:tab-parents}. Counting the allowed pairs $(\kappa,\omega)$ in each case gives $52+4+36+36=128$ formal triples. Comparing with \cref{scf:tab-alphabet}, the triples outside $\mathbb O$ are of two kinds: $\theta=\mathsf c$, $\omega=\mathsf D$ and $\kappa$ a line with $\gamma\ne\mathsf r_3$; and $\theta=\mathsf d$, $\omega=\mathsf C$ and $\kappa$ a line with $\alpha\ne\mathsf r_2$. There are $12$ of each.

Let $t$ be of the first kind, $t=(\mathsf c,\langle\alpha,\gamma\rangle,\mathsf D)$ with $\gamma\ne\mathsf r_3$, and suppose $w$ extends below. Then $w(1,0)$ is the line $\langle\mathsf g_0,\gamma\rangle$. By the thin matching rules (\cref{app:tables}), the tile at $(1,-1)$ has thin type $\mathsf b$, so it is a corner; its north endpoint is $\mathsf g_0$, so it is $\angle\mathsf d$, with west endpoint $\mathsf r_3$. The tile at $(0,-1)$ has thin type $\mathsf d$, so it is a line with black $\mathsf D$ or $\mathsf E$. Its north bit equals the south bit of $\mathsf D$, which is $1$, so its black is $\mathsf D$. Its east endpoint equals the west endpoint $\mathsf r_3$ of $\angle\mathsf d$, so its horizontal atom is $\mathsf r_3$. But thin $\mathsf d$ with black $\mathsf D$ forbids $\gamma=\mathsf r_3$. So the extension does not exist.

Let $t$ be of the second kind, $t=(\mathsf d,\langle\alpha,\gamma\rangle,\mathsf C)$ with $\alpha\ne\mathsf r_2$, and suppose $w$ extends to the left. Then $w(0,1)$ is the line $\langle\alpha,\mathsf g_1\rangle$. The tile at $(-1,1)$ has thin type $\mathsf b$, so it is a corner; its east endpoint is $\mathsf g_1$, so it is $\angle\mathsf d$, with south endpoint $\mathsf r_2$. The tile at $(-1,0)$ has thin type $\mathsf c$, so it is a line with black $\mathsf B$ or $\mathsf C$. Its east bit equals the west bit of $\mathsf C$, which is $1$, so its black is $\mathsf C$. Its north endpoint equals $\mathsf r_2$, so its vertical atom is $\mathsf r_2$. But thin $\mathsf c$ with black $\mathsf C$ forbids $\alpha=\mathsf r_2$. So this extension does not exist either.
\end{proof}

\begin{proposition}[Every tiling has a parent]\label{scf:base-desub}
$\mathbb O$ tiles. For every tiling $f$ of $\mathbb O$ there are a unique coset $c+2\Z^2$ and, for each representative $c$, a unique tiling $z$ of $\mathbb O$ with $f(c+2v+\varsigma)=\varphi(z(v))(\varsigma)$ for all $v\in\Z^2$ and $\varsigma\in[0,2)^2$.
\end{proposition}

\begin{proof}
Let $c+2\Z^2$ be the coset of \cref{scf:thin}; no other coset has thin type $\mathsf a$ at all its sites, so the coset is unique. Fix a representative $c$. For each $v$, the block $f|_{c+2v+[0,2)^2}$ has the thin pattern of \cref{scf:parents}, and the tiling supplies the four cells below and to the left. So it equals $\varphi(z(v))$ for a unique $z(v)\in\mathbb O$. By \cref{scf:base}(b), $z$ is legal, since $\varphi(z)$ is a translate of $f$. For existence, $\varphi^m(t)$ is a legal pattern on $[0,2^m)^2$ for every $t\in\mathbb O$ and $m\ge0$, by \cref{scf:base}(a),(b). So every square carries a legal pattern, and $\mathbb O$ tiles by K\"onig's lemma, as in the proof of \cref{set:compact}.
\end{proof}

The proposition holds for every tiling. No assumption is made that $f$ lies in the closure of the substituted patterns, so tilings with faults along lines or quadrants are covered.

\subsection{The ruler decoration}\label{scf:sub-ruler}

We add a column bit and a row bit to every tile. A \emph{red crossing} is a line $\langle\alpha,\gamma\rangle$ with both atoms red, that is $\alpha\in\{\mathsf r_0,\mathsf r_2\}$ and $\gamma\in\{\mathsf r_1,\mathsf r_3\}$.

\begin{definition}[The scaffold]\label{scf:def-ruler}
The scaffold $\mathbb A$ is the set of tuples $\mathfrak a=(\theta,\kappa,\omega,b_1,b_2)$ with $(\theta,\kappa,\omega)\in\mathbb O$ and $b_1,b_2\in\{0,1\}$ that satisfy:
\begin{itemize}
\item[(R1)] if $\theta\in\{\mathsf b,\mathsf c\}$ then $b_1=0$;
\item[(R2)] if $\theta\in\{\mathsf b,\mathsf d\}$ then $b_2=0$;
\item[(R3)] if $\kappa$ is a corner then $b_1=b_2$;
\item[(R4)] if $\kappa$ is a red crossing then $b_1\ne b_2$.
\end{itemize}
The \emph{column bit} is $\chi_1(\mathfrak a)=b_1$ and the \emph{row bit} is $\chi_2(\mathfrak a)=b_2$. The north and south colours of $\mathfrak a$ are the colours of $(\theta,\kappa,\omega)$ in $\mathbb O$ together with $b_1$; the east and west colours are those colours together with $b_2$. The \emph{base} of $\mathfrak a$ is $\bar{\mathfrak a}=(\theta,\kappa,\omega)$.
\end{definition}

So $\mathbb A$ is a Wang set. \Cref{scf:tab-letters} counts it: $184+4+44+44=276$ letters. By construction, two letters that match in direction $\mathbf e_2$ have the same column bit, and two letters that match in direction $\mathbf e_1$ have the same row bit. We call this the \emph{edgewise bit property}. It holds for every matching pair, not only for pairs that occur in tilings, and the certificates of \cref{scf:sub-parser} use it in that form.

\begin{definition}[Decorated substitution]\label{scf:def-phiprime}
For $\mathfrak a\in\mathbb A$ with bits $(b_1,b_2)$, the pattern $\varphi'(\mathfrak a)$ on $[0,2)^2$ has base $\varphi(\bar{\mathfrak a})$ and the bits of \cref{scf:tab-subst} (right). Patterns $\varphi'(w)$ are defined blockwise as in \cref{scf:def-phi}. Put $\sigma=\varphi'\circ\varphi'$, that is, $\sigma(\mathfrak a)(2\tau+\varsigma)=\varphi'\bigl(\varphi'(\mathfrak a)(\tau)\bigr)(\varsigma)$ for $\tau,\varsigma\in[0,2)^2$. Then $\sigma(\mathfrak a)$ is a pattern on $[0,4)^2$, and $\sigma^n(\mathfrak a)$ is a pattern on $[0,4^n)^2$ with $\sigma^{n+1}(\mathfrak a)(4^n d+y)=\sigma^n\bigl(\sigma(\mathfrak a)(d)\bigr)(y)$ for $d\in[0,4)^2$ and $y\in[0,4^n)^2$.
\end{definition}

\begin{lemma}[Decorated substitution is sound]\label{scf:ruler-forward}
\begin{itemize}
\item[(a)] For every $\mathfrak a\in\mathbb A$, $\varphi'(\mathfrak a)$ is a legal pattern of letters of $\mathbb A$.
\item[(b)] For every pattern $w$ of $\mathbb A$ on a set $U\subseteq\Z^2$, $\varphi'(w)$ is legal if and only if $w$ is legal.
\item[(c)] $\varphi'$ is injective.
\end{itemize}
\end{lemma}

\begin{proof}
(a) By \cref{scf:base}(a) the four bases lie in $\mathbb O$ and match. Let $\mathfrak a=(\theta,\kappa,\omega,b_1,b_2)$. Child $(0,0)$ has thin type $\mathsf a$, so (R1) and (R2) are void; it has thick type $\kappa$ and bits $(1-b_1,1-b_2)$, and (R3), (R4) are unchanged when both bits are complemented. Child $(1,1)$ has thin type $\mathsf b$, a corner and bits $(0,0)$, which satisfies (R1)--(R3). Child $(1,0)$ has thin type $\mathsf c$ and $b_1=0$, which gives (R1). Its thick type $\langle\partial_{\mathrm S}\theta,\partial_{\mathrm E}\kappa\rangle$ is a red crossing only if $\partial_{\mathrm S}\theta$ is red, that is $\theta\in\{\mathsf b,\mathsf d\}$. Then (R2) for $\mathfrak a$ gives $b_2=0$, so the child has bits $(0,1)$ and satisfies (R4). Child $(0,1)$ is symmetric: thin type $\mathsf d$ and row bit $0$; its thick type $\langle\partial_{\mathrm N}\kappa,\partial_{\mathrm W}\theta\rangle$ is a red crossing only if $\theta\in\{\mathsf b,\mathsf c\}$, and then $b_1=0$ and the child has bits $(1,0)$. Inside the block, the vertical edges carry column bits $1-b_1,1-b_1$ and $0,0$, and the horizontal edges carry row bits $1-b_2,1-b_2$ and $0,0$.

(b) The bases match or fail as in \cref{scf:base}(b). Across a vertical seam between the blocks of $\mathfrak a$ (left) and $\mathfrak a'$ (right), the lower fine edge carries row bits $1-b_2$ and $1-b_2'$, and the upper one carries $0$ and $0$. Across a horizontal seam, the left fine edge carries column bits $1-b_1$ and $1-b_1'$, and the right one carries $0$ and $0$. So the fine bits match if and only if the coarse bits match.

(c) Child $(0,0)$ shows $\kappa$, $\omega$ and the complements of both bits, and child $(1,1)$ shows $\theta$.
\end{proof}

\begin{lemma}[Every decorated tiling has a parent]\label{scf:ruler-desub}
For every tiling $f$ of $\mathbb A$ there are a unique coset $c+2\Z^2$ and, for each representative $c$, a unique tiling $z$ of $\mathbb A$ with $f(c+2v+\varsigma)=\varphi'(z(v))(\varsigma)$ for all $v\in\Z^2$ and $\varsigma\in[0,2)^2$.
\end{lemma}

\begin{proof}
The bases of $f$ form a tiling $\bar f$ of $\mathbb O$. Let $c$ and $\bar z$ be given by \cref{scf:base-desub} for $\bar f$. By the edgewise bit property, $\chi_1(f(x))$ depends only on the column $x_1$ and $\chi_2(f(x))$ only on the row $x_2$; write them $X(x_1)$ and $Y(x_2)$. A column $c_1+2v_1+1$ has thin types $\mathsf c$ and $\mathsf b$, so $X=0$ there by (R1). A row $c_2+2v_2+1$ has thin types $\mathsf d$ and $\mathsf b$, so $Y=0$ there by (R2). Define
\[
z(v)=\bigl(\bar z(v),\ 1-X(c_1+2v_1),\ 1-Y(c_2+2v_2)\bigr).
\]
Comparing with \cref{scf:tab-subst}, $f(c+2v+\varsigma)=\varphi'(z(v))(\varsigma)$ for all $v$ and $\varsigma$. The bits of $z$ are constant on columns and on rows and the bases form the tiling $\bar z$, so $z$ is legal once each $z(v)$ lies in $\mathbb A$.

Conditions (R3) and (R4) for $z(v)$ follow from the same conditions for $f(c+2v)$, which has the thick type of $\bar z(v)$ and both bits complemented. For (R1) and (R2) we go down two levels. Apply \cref{scf:base-desub} to $\bar z$: there are $c'$ and a tiling $\bar w$ of $\mathbb O$ with $\bar z(c'+2m+\tau)=\varphi(\bar w(m))(\tau)$. Put $o=c+2c'$. Then $\bar f(o+4m+y)=\varphi^2(\bar w(m))(y)$ for $y\in[0,4)^2$. In this depth-two block, \cref{scf:tab-witness} shows that the cells $(2,1)$ and $(2,3)$ have thin type $\mathsf d$ and that exactly one of them is a red crossing, whatever the thin type of $\bar w(m)$. The row bits of both cells are $0$ by (R2). By (R4) the red crossing has column bit $1$, so $X(o_1+4m_1+2)=1$. In the same way the cells $(1,2)$ and $(3,2)$ have thin type $\mathsf c$, hence column bit $0$, and one of them is a red crossing; so $Y(o_2+4m_2+2)=1$.

Now let $v=c'+2m+\tau$. The letter $z(v)$ has thin type $\mathsf b$ or $\mathsf c$ exactly when $\tau_1=1$. Then $c_1+2v_1=o_1+4m_1+2$, so $\chi_1(z(v))=1-X(o_1+4m_1+2)=0$, which is (R1). In the same way, thin type $\mathsf b$ or $\mathsf d$ means $\tau_2=1$, and then $\chi_2(z(v))=0$, which is (R2).

Uniqueness: the coset is the thin-$\mathsf a$ coset of $\bar f$, the bases of $z$ are unique by \cref{scf:base-desub}, and the bits of $z$ are read off from $f$ at child $(0,0)$.
\end{proof}

One level does not suffice for (R1) and (R2). The witnesses for a coarse column of thin type $\mathsf b$ or $\mathsf c$ sit in the depth-two block, and the offset $o=c+2c'$ combines both levels.

\begin{proposition}[The scaffold]\label{prop:scaffold}
The scaffold $\mathbb A$ is a Wang set of $276$ letters, and it tiles. Its substitution $\sigma=(\varphi')^2$ has expansion $4$ and the following properties.
\begin{itemize}
\item[(i)] $\sigma$ is injective, and $\sigma^n(\mathfrak a)$ is a legal pattern on $[0,4^n)^2$ for every $\mathfrak a\in\mathbb A$ and $n\ge0$.
\item[(ii)] For every pattern $w$ of $\mathbb A$ on a set $U\subseteq\Z^2$, $\sigma(w)$ is legal if and only if $w$ is legal.
\item[(iii)] For every tiling $f$ of $\mathbb A$ there is a unique class $\pi(f)\in(\Z/4\Z)^2$, the \emph{phase} of $f$, such that for every representative $p$ of $\pi(f)$ there is a unique tiling $y$ of $\mathbb A$ with $f(p+4v+i)=\sigma(y(v))(i)$ for all $v\in\Z^2$ and $i\in[0,4)^2$.
\item[(iv)] For every $\mathfrak a\in\mathbb A$ and $i=(i_1,i_2)\in[0,4)^2$,
\[
\chi_1\bigl(\sigma(\mathfrak a)(i)\bigr)=\begin{cases}\chi_1(\mathfrak a)&i_1=0,\\0&i_1\in\{1,3\},\\1&i_1=2,\end{cases}\qquad
\chi_2\bigl(\sigma(\mathfrak a)(i)\bigr)=\begin{cases}\chi_2(\mathfrak a)&i_2=0,\\0&i_2\in\{1,3\},\\1&i_2=2.\end{cases}
\]
\item[(v)] Letters that match in direction $\mathbf e_2$ have equal column bits, and letters that match in direction $\mathbf e_1$ have equal row bits.
\end{itemize}
\end{proposition}

\begin{proof}
(i) and (ii) follow from \cref{scf:ruler-forward} applied twice, and (v) is the edgewise bit property. By (i), every square carries a legal pattern, so $\mathbb A$ tiles by K\"onig's lemma.

(iii) Let $f$ be a tiling. \Cref{scf:ruler-desub} gives $f(c+2v+\varsigma)=\varphi'(z(v))(\varsigma)$, and again $z(c'+2m+\tau)=\varphi'(y(m))(\tau)$. Substituting, $f(c+2c'+4m+2\tau+\varsigma)=\sigma(y(m))(2\tau+\varsigma)$, so the class of $p=c+2c'$ works; for another representative $p+4m_0$ use $y(\cdot+m_0)$. Conversely, suppose $f(p+4m+i)=\sigma(y(m))(i)$ for a tiling $y$. Then $f(p+2v+\varsigma)=\varphi'(z'(v))(\varsigma)$ with $z'(2m+\tau)=\varphi'(y(m))(\tau)$, and $z'$ is a tiling by \cref{scf:ruler-forward}. By the uniqueness in \cref{scf:ruler-desub}, $p\equiv c\pmod2$. Write $p=c+2d$; uniqueness gives $z'=z(\cdot+d)$. So $z(d+2m+\tau)=\varphi'(y(m))(\tau)$, and uniqueness for $z$ gives $d\equiv c'\pmod 2$. Hence $p\equiv c+2c'\pmod4$. For a fixed representative, $y$ is unique because $\sigma$ is injective.

(iv) Write $i_1=2\tau_1+\varsigma_1$. By \cref{scf:tab-subst}, the column bit of $\varphi'(\mathfrak b)(\varsigma)$ is $1-\chi_1(\mathfrak b)$ if $\varsigma_1=0$ and $0$ if $\varsigma_1=1$. Applying this twice: for $i_1=0$ we get $1-(1-\chi_1(\mathfrak a))=\chi_1(\mathfrak a)$; for $i_1\in\{1,3\}$ we get $0$; for $i_1=2$ the middle letter has column bit $0$, so we get $1$. The row bit is the same with $\tau_2,\varsigma_2$.
\end{proof}

Iterating (iv), the letter $\sigma^n(\mathfrak a)(0,0)$ has the bits of $\mathfrak a$ for every $n$. In a tiling of phase $\pi$, the column bit of a column $x_1$ is $0$ if $x_1-p_1\equiv1,3\pmod4$ and $1$ if $x_1-p_1\equiv2\pmod 4$, for any representative $p$ of $\pi$. Only the columns $x_1\equiv p_1\pmod4$ inherit a bit from the parent. So the column bits form a ruler sequence, like the parity of the $2$-adic valuation of a column index. The rows behave the same way.

\subsection{Reading a block from a fixed number of sites}\label{scf:sub-ports}

A depth-$n$ block $\sigma^n(\mathfrak a)$ has $16^n$ cells, but its root letter can be read from at most $\ell$ of them, at every depth.

Let $\mathcal P$ be the set of unordered pairs of distinct letters of $\mathbb A$; it has $\binom{276}{2}=37950$ elements. For each $p=\{\mathfrak a,\mathfrak b\}\in\mathcal P$ fix once and for all a child $\delta(p)\in[0,4)^2$ with $\sigma(\mathfrak a)(\delta(p))\ne\sigma(\mathfrak b)(\delta(p))$. It exists because $\sigma$ is injective. Put $\nu(p)=\{\sigma(\mathfrak a)(\delta(p)),\sigma(\mathfrak b)(\delta(p))\}$, again an element of $\mathcal P$. Define offsets $y_n(p)\in[0,4^n)^2$ by
\[
y_0(p)=(0,0),\qquad y_{n+1}(p)=4^n\delta(p)+y_n(\nu(p)),
\]
and put
\[
P_n=\{(0,0)\}\cup\{y_n(p):p\in\mathcal P\}\subseteq[0,4^n)^2,\qquad \operatorname{Port}_n(u)=u+P_n,\qquad \ell=\tbinom{276}{2}+1=37951 .
\]
So $|\operatorname{Port}_n(u)|\le\ell$; coincident offsets are counted once. The \emph{readout} of $\mathfrak a$ at depth $n$ is the restriction of $\sigma^n(\mathfrak a)$ to $P_n$. A map $g$ into $\mathbb A$ defined on $\operatorname{Port}_n(v)$ \emph{reads} $\mathfrak a$ if $g(v+y)=\sigma^n(\mathfrak a)(y)$ for every $y\in P_n$.

\begin{lemma}[Pair ports]\label{scf:ports}
\begin{itemize}
\item[(a)] For every $n\ge0$, the readout at depth $n$ is injective on $\mathbb A$.
\item[(b)] Every $y\in P_{n+1}$ has the form $y=4^nd+y'$ with $d\in[0,4)^2$ and $y'\in P_n$, and then $\sigma^{n+1}(\mathfrak a)(y)=\sigma^n\bigl(\sigma(\mathfrak a)(d)\bigr)(y')$ for every $\mathfrak a$.
\item[(c)] $(0,0)\in P_n$, and $\sigma^n(\mathfrak a)(0,0)$ has the column bit and the row bit of $\mathfrak a$.
\end{itemize}
\end{lemma}

\begin{proof}
(a) We show $\sigma^n(\mathfrak a)(y_n(p))\ne\sigma^n(\mathfrak b)(y_n(p))$ for every $p=\{\mathfrak a,\mathfrak b\}$ and every $n$. For $n=0$ the two values are $\mathfrak a$ and $\mathfrak b$. For $n+1$, \cref{scf:def-phiprime} gives $\sigma^{n+1}(\mathfrak a)(y_{n+1}(p))=\sigma^n(\mathfrak a')(y_n(\nu(p)))$ with $\mathfrak a'=\sigma(\mathfrak a)(\delta(p))$, and the same for $\mathfrak b$ with $\mathfrak b'=\sigma(\mathfrak b)(\delta(p))$. Since $\nu(p)=\{\mathfrak a',\mathfrak b'\}$, the induction hypothesis applies. So the offset $y_n(\{\mathfrak a,\mathfrak b\})$ separates any two distinct letters.

(b) The offset $(0,0)$ is $4^n\cdot(0,0)+(0,0)$, and $y_{n+1}(p)=4^n\delta(p)+y_n(\nu(p))$ with $y_n(\nu(p))\in P_n$. The identity is \cref{scf:def-phiprime}.

(c) This is \cref{prop:scaffold}(iv) at $i=(0,0)$, applied $n$ times.
\end{proof}

Part (b) says that each port of a depth-$(n+1)$ block is a port of one of its $16$ children. So the readouts of all depths are nested and their number of sites does not grow with $n$. The same construction works for any injective substitution on $m$ letters, with $\binom m2+1$ sites.

\subsection{Anchors and certificates}\label{scf:sub-parser}

For $u\in\Z^2$ and $n\ge0$ put
\[
\operatorname{Anc}_n(u)=\{u,\ u+4^n\mathbf e_1,\ u+4^n\mathbf e_2\}.
\]

\begin{definition}[Anchors]\label{scf:def-anchor}
A map $g$ into $\mathbb A$ whose domain contains $\operatorname{Anc}_n(u)$ is a \emph{good $n$-anchor at $u$} if $\chi_1(g(u))\ne\chi_1(g(u+4^n\mathbf e_1))$ and $\chi_2(g(u))\ne\chi_2(g(u+4^n\mathbf e_2))$.
\end{definition}

The condition involves three sites and their labels only. It does not assume that the labels come from a block of depth $n$; the certificates below derive this.

\begin{definition}[Certificates]\label{scf:def-cert}
A \emph{certificate} $\Gamma$ consists of a finite vertex set $V$, a map $\rho\colon V\to\Z^2$ (its \emph{grid map}), a set of \emph{edges}, and a finite set $\partial\Gamma\subseteq\Z^2$ of sites, its \emph{boundary}. Each edge is an ordered pair $(v,v')$ of vertices with a label $e\in\{\mathbf e_1,\mathbf e_2\}$ such that $\rho(v')=\rho(v)+e$. Each boundary site $x$ has exactly one \emph{boundary vertex} $v_x$, with $\rho(v_x)=x$. A \emph{homomorphism} of $\Gamma$ is a map $\psi\colon V\to\mathbb A$ such that $\psi(v)$ and $\psi(v')$ match in direction $e$ for every edge $(v,v')$ with label $e$. For a boundary site $x$ we write $\psi(x)$ for $\psi(v_x)$.

A \emph{rooted tree decomposition} of $\Gamma$ is a finite rooted tree with a \emph{bag} $\beta_t\subseteq V$ at each node $t$ such that every vertex lies in some bag, the nodes whose bags contain a given vertex form a connected subtree, both ends of every edge lie in a common bag, and the root bag contains every boundary vertex. Its \emph{width} is the largest bag size.
\end{definition}

Distinct vertices may have the same image under $\rho$, and a homomorphism may give them different letters. Only the edges and the boundary tie vertices together. This freedom is what keeps the width bounded: a certificate may use the same site in many independent sub-certificates without identifying the copies. Certificates of bounded tree width are the structures for which the existential pebble game is decisive, by the work of Dalmau, Kolaitis and Vardi \cite{DKV02}; \cref{lem:transfer} below is the form we need.

\paragraph{A finite recognition chart.} The next lemma replaces an explicit table of local phases by compactness. We need only that the constants exist; they are never computed.

\begin{lemma}[Recognition chart]\label{scf:chart}
\begin{itemize}
\item[(a)] For every finite set $D\subseteq\Z^2$ there is $L(D)\ge0$ such that, for every $u\in\Z^2$, every legal pattern of $\mathbb A$ on $u+D+[-L(D),L(D)]^2$ restricts to an extendable pattern on $u+D$.
\item[(b)] There is $L_0\ge8$ such that any two tilings of $\mathbb A$ that agree on $[-L_0,L_0]^2$ have the same phase.
\end{itemize}
\end{lemma}

\begin{proof}
(a) Legality and extendability are invariant under translation, so take $u=0$. Let $w$ be a pattern on $D$ that is not extendable. Suppose that for every $L$ it extends to a legal pattern on $D+[-L,L]^2$. These extensions form a finitely branching tree under restriction, so K\"onig's lemma gives a legal pattern on $\Z^2$ that extends $w$, a contradiction. So $w$ has a radius $L_w$ beyond which it has no legal extension. There are finitely many patterns on $D$; let $L(D)$ be the largest $L_w$.

(b) For a class $\pi\in(\Z/4\Z)^2$, the set of tilings of phase $\pi$ is closed in $\mathbb A^{\Z^2}$. Indeed, let $f_k\to f$ be tilings of phase $\pi$ and fix a representative $p$. Write $f_k(p+4v+i)=\sigma(y_k(v))(i)$. Since $\sigma$ is injective, $y_k(v)$ is determined by $f_k$ on the block $p+4v+[0,4)^2$, so $y_k$ converges to a map $y$ with $f(p+4v+i)=\sigma(y(v))(i)$. The pattern $\sigma(y)$ is a translate of $f$, hence legal, so $y$ is a tiling by \cref{prop:scaffold}(ii). Now suppose (b) fails. Then for every $L$ there are tilings $f_L,f'_L$ that agree on $[-L,L]^2$ and have different phases. Passing to a subsequence, the phases are constant classes $\pi\ne\pi'$, and $f_L\to f$, $f'_L\to f'$ by compactness of $\mathbb A^{\Z^2}$. The limits agree on every box, so $f=f'$, and by closedness $f$ has both phases $\pi$ and $\pi'$. This contradicts \cref{prop:scaffold}(iii). Enlarging $L_0$ keeps (b) true, so we may take $L_0\ge8$.
\end{proof}

The phase of a translate changes by the translation: if $f(p+4v+i)=\sigma(y(v))(i)$ then $f(\cdot+u)$ has phase $\pi(f)-u$. So by (b), two tilings that agree on $u+[-L_0,L_0]^2$ have the same phase, for every $u$.

\begin{definition}[Constants]\label{scf:def-constants}
Fix $L_0$ as in \cref{scf:chart}(b). Put $\mathsf D_0=[-L_0,L_0]^2$, $Q=[-L_0-1,L_0+1]^2$, $L_1=L(Q)$ and $\varrho=L_0+1+L_1$. Put $L_2=L([0,5)^2)$ and $L_3=L([0,32)^2)$. Define
\[
B=\max\Bigl\{\ell(2\varrho+3)^2+2\ell+26,\ \ (2L_2+5)^2,\ \ 3+\ell\,(2L_3+32)^2\Bigr\}.
\]
\end{definition}

The integer $B$ depends only on $\mathbb A$, $\sigma$ and the fixed choice of $\delta$. It does not depend on any depth, site, distance or machine. Its first term bounds the bags of the certificates; the other two are the sizes of the fixed boxes used in \cref{prop:anchors}. In particular $B\ge\ell+4$. We only use that $B$ exists.

\paragraph{Corridors.} A \emph{corridor walk} is a sequence $w_0,\dots,w_m\in\Z^2$ in which each step $w_{k+1}-w_k$ is one of $\pm\mathbf e_1,\pm\mathbf e_2$ and each coordinate of $w_k$ is monotone in $k$. Its \emph{corridor} is $S=\bigcup_k\bigl(w_k+[-\varrho,\varrho]^2\bigr)$.

\begin{lemma}[Phase transport]\label{scf:corridor}
Let $w_0,\dots,w_m$ be a corridor walk with corridor $S$, and let $\mathfrak c$ be a legal pattern of $\mathbb A$ on $S$. Suppose that $\chi_1(\mathfrak c(w_0))\ne\chi_1(\mathfrak c(w_0+4\mathbf e_1))$ and $\chi_2(\mathfrak c(w_0))\ne\chi_2(\mathfrak c(w_0+4\mathbf e_2))$. Let $v\equiv w_0\pmod4$ with $v+[0,8)^2\subseteq w_k+\mathsf D_0$ for some $k$. Then there are letters $\mathfrak a,\mathfrak b_1,\mathfrak b_2$ such that, for $i\in[0,4)^2$ and $e\in\{\mathbf e_1,\mathbf e_2\}$,
\[
\mathfrak c(v+i)=\sigma(\mathfrak a)(i),\qquad \mathfrak c(v+4e+i)=\sigma(\mathfrak b_e)(i),
\]
and $\mathfrak a$ and $\mathfrak b_e$ match in direction $e$.
\end{lemma}

\begin{proof}
For every $k$ the pattern $\mathfrak c$ is legal on $w_k+Q+[-L_1,L_1]^2=w_k+[-\varrho,\varrho]^2$, so by \cref{scf:chart}(a) its restriction to $w_k+Q$ is extendable. Let $g_k$ be a tiling that agrees with $\mathfrak c$ on $w_k+Q$, and let $\pi_k$ be its phase. By the remark after \cref{scf:chart}, $\pi_k$ depends only on $\mathfrak c|_{w_k+\mathsf D_0}$. Since $w_{k+1}+\mathsf D_0\subseteq w_k+Q$, the tiling $g_k$ also agrees with $\mathfrak c$ on $w_{k+1}+\mathsf D_0$; hence $\pi_{k+1}=\pi_k$. So all $\pi_k$ equal $\pi_0$.

The sites $w_0$, $w_0+4\mathbf e_1$ and $w_0+4\mathbf e_2$ lie in $w_0+\mathsf D_0$ because $L_0\ge8$. Let $p$ represent $\pi_0$. If $w_{0,1}\not\equiv p_1\pmod4$, then the columns $w_{0,1}$ and $w_{0,1}+4$ have the same nonzero residue modulo $4$ relative to $p_1$, so by \cref{prop:scaffold}(iii),(iv) they have the same column bit in $g_0$; this contradicts the hypothesis. So $w_{0,1}\equiv p_1$, and in the same way $w_{0,2}\equiv p_2$, using rows and row bits. Thus $\pi_k$ is the class of $w_0$ for every $k$, and also the class of $v$.

Take $k$ with $v+[0,8)^2\subseteq w_k+\mathsf D_0$ and use the representative $v$: by \cref{prop:scaffold}(iii) there is a tiling $y$ with $g_k(v+4m+i)=\sigma(y(m))(i)$. Put $\mathfrak a=y(0)$ and $\mathfrak b_e=y(e)$. The cells $v+i$ and $v+4e+i$ lie in $v+[0,8)^2$, where $g_k$ agrees with $\mathfrak c$. Since $y$ is a tiling, $\mathfrak a$ and $\mathfrak b_e$ match in direction $e$.
\end{proof}

The walk carries the phase step by step through overlapping windows. This is the only way the phase fixed at the anchor reaches a distant block, and every certificate below contains such a walk.

\paragraph{The parser.} For $n\ge0$ and $u,z\in\Z^2$, the \emph{block} of $z$ at depth $n$ relative to $u$ is $u+4^nz+[0,4^n)^2$, and its ports are $\operatorname{Port}_n(u+4^nz)$.

\begin{theorem}[Parser]\label{thm:parser}
Let $B$ be as in \cref{scf:def-constants}. For all $n\ge0$, $u,z\in\Z^2$, $e\in\{\mathbf e_1,\mathbf e_2\}$ and $s\in u+4^nz+[0,4^n)^2$ there are three certificates:
\begin{itemize}
\item the \emph{type certificate} $\Gamma^{\mathrm{ty}}_n(u,z)$, with boundary $\operatorname{Anc}_n(u)\cup\operatorname{Port}_n(u+4^nz)$;
\item the \emph{edge certificate} $\Gamma^{\mathrm{ed}}_n(u,z,e)$, with boundary $\operatorname{Anc}_n(u)\cup\operatorname{Port}_n(u+4^nz)\cup\operatorname{Port}_n(u+4^n(z+e))$;
\item the \emph{point certificate} $\Gamma^{\mathrm{pt}}_n(u,z,s)$, with boundary $\operatorname{Anc}_n(u)\cup\operatorname{Port}_n(u+4^nz)\cup\{s\}$.
\end{itemize}
Each has a rooted tree decomposition of width at most $B$. Let $\psi$ be a homomorphism of one of them whose values on $\operatorname{Anc}_n(u)$ form a good $n$-anchor at $u$. Then:
\begin{itemize}
\item[(T)] there is $\mathfrak a\in\mathbb A$ such that $\psi$ reads $\mathfrak a$ on $\operatorname{Port}_n(u+4^nz)$;
\item[(E)] for the edge certificate, moreover $\psi$ reads some $\mathfrak b$ on $\operatorname{Port}_n(u+4^n(z+e))$, and $\mathfrak a$ and $\mathfrak b$ match in direction $e$;
\item[(P)] for the point certificate, moreover $\psi(s)=\sigma^n(\mathfrak a)(s-u-4^nz)$.
\end{itemize}
The bound $B$ does not depend on $n$, $u$, $z$, $e$, $s$, on the distance from the anchor to the block, or on any machine.
\end{theorem}

By \cref{scf:ports}(a) the letters $\mathfrak a$ and $\mathfrak b$ are unique. The query block may be arbitrarily far from the anchor.

\begin{proof}
We build the three certificates together, by induction on $n$, and check the decomposition and the conclusions at each step.

\emph{Depth $0$.} Here $P_0=\{(0,0)\}$, so the ports of $z$ form the single site $u+z$, and $s=u+z$. The type and point certificates have one vertex for each boundary site and no edges. The edge certificate adds one edge from $u+z$ to $u+z+e$ with label $e$. One bag holds everything; it has at most $5$ vertices. Conclusions (T) and (P) hold with $\mathfrak a=\psi(u+z)$, and (E) holds because of the edge. No anchor condition is needed at depth $0$.

\emph{From depth $n$ to $n+1$: the skeleton.} Fix $u$, $z$, and $e$ or $s$ as required. Coarse sites $x\in\Z^2$ stand for the depth-$n$ blocks $u+4^nx+[0,4^n)^2$; these blocks tile the plane. The depth-$(n+1)$ block of $z$ is the union of the depth-$n$ blocks of $4z+d$, $d\in[0,4)^2$. Choose a corridor walk from $w_0=(0,0)$ to $w_m=4z$, say first along $\mathbf e_1$ and then along $\mathbf e_2$, and let $S$ be its corridor. The \emph{skeleton} has one vertex for each site of
\[
\bigcup_{x\in S}\operatorname{Port}_n(u+4^nx),
\]
and, for the point certificate, one more vertex for $s$ if $s$ is not already among these sites. Its grid map is the identity on sites. The skeleton contains the $25$ \emph{midpoint} sites $u+4^n(\iota,\kappa_0)$, $\iota,\kappa_0\in[0,4]$, because each is the corner port of a coarse site of $[0,4]^2\subseteq S$. Write $\mathcal M_0$ for this set. The boundary lies in the skeleton: $\operatorname{Anc}_{n+1}(u)\subseteq\mathcal M_0$, and $\operatorname{Port}_{n+1}(u+4^{n+1}z)$ is contained in the union of the ports of the coarse sites $4z+d$ by \cref{scf:ports}(b). The same holds for $z+e$ in the edge certificate, since $4z+4e+[0,4)^2\subseteq 4z+\mathsf D_0\subseteq S$.

\emph{Bit paths.} For $\iota\in[0,4]$ and $\kappa_0\in[0,3]$ attach a path of fresh vertices along the column of $u+4^n(\iota,\kappa_0)$ from that site to $u+4^n(\iota,\kappa_0+1)$, with edges labelled $\mathbf e_2$; its two end vertices are the skeleton vertices of these midpoints. For $\kappa_0\in[0,4]$ and $\iota\in[0,3]$ attach in the same way a path along the row from $u+4^n(\iota,\kappa_0)$ to $u+4^n(\iota+1,\kappa_0)$, with edges labelled $\mathbf e_1$. There are $40$ paths.

\emph{The sixteen candidates.} For each $c\in[0,4)^2$ put $v_c=u+4^nc$, so that $\operatorname{Anc}_n(v_c)\subseteq\mathcal M_0$. For each $c$ attach a fresh copy of each of the following depth-$n$ certificates:
\begin{itemize}
\item $\Gamma^{\mathrm{ty}}_n(v_c,x-c)$ for every $x\in S$;
\item $\Gamma^{\mathrm{ed}}_n(v_c,x-c,e')$ for every $x$ and $e'\in\{\mathbf e_1,\mathbf e_2\}$ with $x,x+e'\in S$;
\item for the point certificate only, $\Gamma^{\mathrm{pt}}_n(v_c,4z+d_s-c,s)$, where $d_s\in[0,4)^2$ is the child whose block contains $s$.
\end{itemize}
The block of $x-c$ relative to $v_c$ is the block of $x$ relative to $u$. So the boundary of each copy consists of skeleton sites, and we identify each boundary vertex of the copy with the skeleton vertex of the same site. All other vertices of the copy are fresh. In particular, the sixteen candidates share the same port vertices: they read the same physical blocks.

\emph{The conclusions.} Let $\psi$ be a homomorphism of the depth-$(n+1)$ certificate whose values on $\operatorname{Anc}_{n+1}(u)$ form a good $(n+1)$-anchor at $u$. Write $m_{\iota\kappa_0}=\psi(u+4^n(\iota,\kappa_0))$.

Step 1: one good candidate. Along a column path every edge joins letters that match in direction $\mathbf e_2$, so all letters on the path have the same column bit (\cref{prop:scaffold}(v)). Hence $\chi_1(m_{\iota\kappa_0})=\chi_1(m_{\iota0})$ for all $\iota,\kappa_0\in[0,4]$, and in the same way $\chi_2(m_{\iota\kappa_0})=\chi_2(m_{0\kappa_0})$. Since $\chi_1(m_{00})\ne\chi_1(m_{40})$, some $\iota\in[0,3]$ has $\chi_1(m_{\iota0})\ne\chi_1(m_{\iota+1,0})$. Since $\chi_2(m_{00})\ne\chi_2(m_{04})$, some $\kappa_0\in[0,3]$ has $\chi_2(m_{0\kappa_0})\ne\chi_2(m_{0,\kappa_0+1})$. Then $\chi_1(m_{\iota\kappa_0})\ne\chi_1(m_{\iota+1,\kappa_0})$ and $\chi_2(m_{\iota\kappa_0})\ne\chi_2(m_{\iota,\kappa_0+1})$. So $\psi$ is a good $n$-anchor at $v_c$ for $c=(\iota,\kappa_0)$. The construction attached all sixteen candidates in advance; which one is good depends on $\psi$, and the argument below uses only that one.

Step 2: a common coarse pattern. The restriction of $\psi$ to a copy attached for $c$ is a homomorphism of that copy, with the same values on its boundary. By the induction hypothesis (T), for every $x\in S$ the map $\psi$ reads some letter on $\operatorname{Port}_n(u+4^nx)$; by \cref{scf:ports}(a) it is unique, and we call it $\mathfrak c(x)$. It depends only on the values of $\psi$ at the skeleton vertices of these ports, so every copy that meets these ports sees the same letter. By (E), $\mathfrak c(x)$ and $\mathfrak c(x+e')$ match in direction $e'$ whenever $x,x+e'\in S$. So $\mathfrak c$ is a legal pattern on $S$.

Step 3: the phase. By \cref{scf:ports}(c), $m_{\iota\kappa_0}=\psi(u+4^n(\iota,\kappa_0))=\sigma^n(\mathfrak c(\iota,\kappa_0))(0,0)$ has the bits of $\mathfrak c(\iota,\kappa_0)$. The good $(n+1)$-anchor therefore gives $\chi_1(\mathfrak c(0,0))\ne\chi_1(\mathfrak c(4,0))$ and $\chi_2(\mathfrak c(0,0))\ne\chi_2(\mathfrak c(0,4))$. Apply \cref{scf:corridor} with $v=4z$, which is congruent to $w_0=(0,0)$ modulo $4$ and satisfies $4z+[0,8)^2\subseteq w_m+\mathsf D_0$ because $L_0\ge8$. It gives a letter $\mathfrak a$ with $\mathfrak c(4z+d)=\sigma(\mathfrak a)(d)$ for $d\in[0,4)^2$, and, for the edge certificate, a letter $\mathfrak b$ with $\mathfrak c(4z+4e+d)=\sigma(\mathfrak b)(d)$, where $\mathfrak a$ and $\mathfrak b$ match in direction $e$.

Step 4: the readout. Let $y\in P_{n+1}$ and write $y=4^nd+y'$ with $y'\in P_n$, by \cref{scf:ports}(b). Then
\[
\begin{aligned}
\psi(u+4^{n+1}z+y)&=\psi\bigl(u+4^n(4z+d)+y'\bigr)=\sigma^n\bigl(\mathfrak c(4z+d)\bigr)(y')\\
&=\sigma^n\bigl(\sigma(\mathfrak a)(d)\bigr)(y')=\sigma^{n+1}(\mathfrak a)(y).
\end{aligned}
\]
So $\psi$ reads $\mathfrak a$ on $\operatorname{Port}_{n+1}(u+4^{n+1}z)$, which is (T). The same computation for $z+e$ gives (E).

Step 5: the point. For the good candidate $c$, the induction hypothesis (P) for the attached copy of $\Gamma^{\mathrm{pt}}_n(v_c,4z+d_s-c,s)$ gives a letter $\mathfrak a''$ read on the ports of $4z+d_s$, with $\psi(s)=\sigma^n(\mathfrak a'')(s-u-4^n(4z+d_s))$. By uniqueness of the readout, $\mathfrak a''=\mathfrak c(4z+d_s)=\sigma(\mathfrak a)(d_s)$. Since $s-u-4^{n+1}z=4^nd_s+\bigl(s-u-4^n(4z+d_s)\bigr)$, \cref{scf:def-phiprime} gives $\psi(s)=\sigma^{n+1}(\mathfrak a)(s-u-4^{n+1}z)$, which is (P).

\emph{The decomposition.} Let $\partial$ be the boundary of the depth-$(n+1)$ certificate. For $k\in[0,m]$ let $\beta_k$ consist of the skeleton vertices at the sites of $\mathcal M_0$, of $\partial$, and of $\operatorname{Port}_n(u+4^nx)$ for every $x\in S\cap\bigl(w_k+[-\varrho-1,\varrho+1]^2\bigr)$. Arrange $\beta_0,\dots,\beta_m$ in a path with root $\beta_0$. A skeleton vertex at a site of $\mathcal M_0\cup\partial$ lies in every bag. Any other skeleton vertex lies in the ports of exactly one coarse site $x$, because the depth-$n$ blocks are disjoint; the set of $k$ with $x\in w_k+[-\varrho-1,\varrho+1]^2$ is an interval, because each coordinate of $w_k$ is monotone in $k$. So the skeleton bags satisfy the connectivity condition. Every $x\in S$ lies in $w_k+[-\varrho,\varrho]^2$ for some $k$, and then its ports and those of every neighbour $x+e'\in S$ lie in $\beta_k$.

Attach the decomposition of each child copy by joining its root bag to a bag that contains the copy's boundary. For a type copy at $x$ or an edge copy at $x,x+e'$, take $\beta_k$ with $x\in w_k+[-\varrho,\varrho]^2$; the anchor sites of the copy lie in $\mathcal M_0$. For the point copy, take $\beta_m$: the coarse site $4z+d_s$ lies in $w_m+[-\varrho,\varrho]^2$, and $s\in\partial$. For a bit path with vertices $x_0,\dots,x_N$, use the bags $\{x_0,x_N,x_k,x_{k+1}\}$, $k<N$, in a path, and join its first bag to $\beta_0$, which contains the midpoints $x_0$ and $x_N$. The fresh vertices of a copy or a path lie only in its own bags, and each shared vertex lies in a connected set of skeleton bags together with the connected set of bags of each copy that contains it, joined at the attachment. Every edge lies inside a copy or a path, where a bag contains it. So this is a rooted tree decomposition, and $\beta_0\supseteq\partial$.

Each $\beta_k$ has at most $\ell(2\varrho+3)^2$ port vertices, $25$ midpoints, and at most $2\ell+1$ further boundary vertices, so $|\beta_k|\le\ell(2\varrho+3)^2+2\ell+26\le B$. The bags of the paths have $4$ vertices, and the bags of the copies have at most $B$ vertices by induction. So the width is at most $B$. The bound is a maximum over finitely many kinds of bags, not a sum over depths, and the length of the corridor changes only the number of bags.
\end{proof}

Two features of the construction carry the proof. First, all sixteen candidates are attached to the same port vertices, so the letter decoded at a coarse site is one letter, whatever certificate looks at it. Second, the point certificate contains the whole corridor from the anchor to the block that holds $s$. A window near the block alone would fix some phase of the decoded pattern, but nothing would tie that phase to the anchor's grid; inside the recursion the child block can lie at coarse distance $3,15,63,\dots$ from the candidate anchors.

The bit paths use the edgewise bit property of \cref{prop:scaffold}(v). If the bits were only known to be constant along the columns and rows of tilings, a bare path would force nothing, since a homomorphism of a certificate need not extend to a tiling. The literal bit colours of \cref{scf:def-ruler} avoid this.

\subsection{Transfer to pebble families, and anchors of every depth}\label{scf:sub-transfer}

We use $E_K$ families for the Wang set $\mathbb A$ as in \cref{set:def-pebble}: nonempty sets of legal patterns of $\mathbb A$ on finite subsets of $\Z^2$, each with at most $K$ sites, closed under restriction and with the extension property for maps on fewer than $K$ sites. By \cref{set:pebble-basic}, a member $h$ extends to a member defined on $\operatorname{dom}h\cup U'$ whenever $|\operatorname{dom}h\cup U'|\le K$.

\begin{lemma}[Transfer]\label{lem:transfer}
Let $H$ be an $E_K$ family for $\mathbb A$. Let $\Gamma$ be a certificate with a rooted tree decomposition of width at most $B'$, and let $\Pi\subseteq\Z^2$ be a finite set of \emph{pinned} sites with $|\Pi|\le K-B'$. Let $h\in H$ have domain $\partial\Gamma\cup\Pi$. Then there is a homomorphism $\psi$ of $\Gamma$ such that $\psi(x)=h(x)$ for every boundary site $x$, and $\psi(v)=h(\rho(v))$ for every vertex $v$ with $\rho(v)\in\Pi$.
\end{lemma}

So a use of a certificate of width at most $B'$ may retain up to $K-B'$ pinned sites beyond its bags, and the certificate then respects the pins wherever it touches them. With $B'=B$ and $K=j=10(B+\ell+100)$, as in \cref{thm:halting}, this allows $9B+10\ell+1000$ pinned sites. The applications in this paper need fewer: \cref{prop:anchors} and \cref{scf:readout} below take $\Pi=\emptyset$, and \cref{sec:refutation} applies certificates only to restrictions of its maps to the certificate boundary, which is again the case $\Pi=\emptyset$. All maps there have at most $\ell+6\le B$ sites, so the single condition $K\ge B$ covers the fixed boxes of \cref{prop:anchors}, the certificates, and the pins and sibling sites of \cref{sec:refutation}.

\begin{proof}
Let $(\beta_t)$ be the decomposition, with root $t_0$; it has width at most $B'$, so $|\rho(\beta_t)|\le B'$. We choose maps $g_t\in H$ with domain $\rho(\beta_t)\cup\Pi$ and $g_t|_\Pi=h|_\Pi$, from the root down. For the root, $\partial\Gamma\subseteq\rho(\beta_{t_0})$, and $h$ extends to $g_{t_0}$ on $\rho(\beta_{t_0})\cup\Pi$, a set of at most $K$ sites. For a child $t$ of a node $t'$, restrict $g_{t'}$ to $\rho(\beta_t\cap\beta_{t'})\cup\Pi$ and extend the restriction to $\rho(\beta_t)\cup\Pi$. Put $\psi(v)=g_t(\rho(v))$ for any $t$ with $v\in\beta_t$. This does not depend on $t$: the nodes whose bags contain $v$ form a subtree, and along each of its edges the two maps agree at $\rho(v)$ by construction. Every edge $(v,v')$ of $\Gamma$ with label $e$ lies in a bag $\beta_t$; then $\psi(v)=g_t(\rho(v))$ and $\psi(v')=g_t(\rho(v)+e)$ are values of one legal pattern at sites $\rho(v)$ and $\rho(v)+e$, so they match in direction $e$. At a boundary site $x$, $\psi(x)=g_{t_0}(x)=h(x)$. If $\rho(v)\in\Pi$, then $\psi(v)=g_t(\rho(v))=h(\rho(v))$.
\end{proof}

This is the standard argument showing that the existential $K$-pebble game decides homomorphism from structures of tree width less than $K$ \cite{KolaitisVardi00,DKV02}. Two vertices with one site in a common bag receive one letter. Vertices with one site in unrelated branches may receive different letters, and \cref{thm:parser} never needs them to agree.

\begin{corollary}[Readout in a family]\label{scf:readout}
Let $H$ be an $E_K$ family for $\mathbb A$ with $K\ge B$. Let $n\ge0$, $u\in\Z^2$, and let $g\in H$ be defined on $\operatorname{Anc}_n(u)\cup\operatorname{Port}_n(u)$ and be a good $n$-anchor at $u$. Then $g$ reads a letter $\mathfrak a$ on $\operatorname{Port}_n(u)$, that is $g(y)=\sigma^n(\mathfrak a)(y-u)$ for every $y\in\operatorname{Port}_n(u)$; the letter is unique. If $g$ is also defined at a site $s\in u+[0,4^n)^2$, then $g(s)=\sigma^n(\mathfrak a)(s-u)$. More generally, for every $z\in\Z^2$: if $g$ is defined on $\operatorname{Anc}_n(u)\cup\operatorname{Port}_n(u+4^nz)$ with a good anchor, it reads a letter on $\operatorname{Port}_n(u+4^nz)$, and the value at any $s$ in the block of $z$ is the corresponding entry of $\sigma^n$ of that letter.
\end{corollary}

\begin{proof}
Restrict $g$ to the boundary of the type certificate $\Gamma^{\mathrm{ty}}_n(u,z)$, or of the point certificate $\Gamma^{\mathrm{pt}}_n(u,z,s)$, and apply \cref{lem:transfer} with $\Pi=\emptyset$ and $B'=B$. The resulting homomorphism has a good anchor, and \cref{thm:parser}(T),(P) give the claims. The letter given by the point certificate is the one read on the ports, by \cref{scf:ports}(a).
\end{proof}

\begin{proposition}[Anchors of every depth]\label{prop:anchors}
Let $H$ be an $E_K$ family for $\mathbb A$ with $K\ge B$. For every $n\ge0$ there are $u\in\Z^2$ and $h\in H$ with domain $\operatorname{Anc}_n(u)$ that is a good $n$-anchor at $u$.
\end{proposition}

\begin{proof}
We induct on $n$. The site $u$ may depend on $H$ and on $n$.

\emph{Depth $0$.} The box $Y=[-L_2,4+L_2]^2$ has $(2L_2+5)^2\le B\le K$ sites, so $H$ has a member $g$ with domain $Y$ (\cref{set:pebble-basic}). It is a legal pattern, so by \cref{scf:chart}(a) its restriction to $[0,5)^2$ is the restriction of a tiling $f$. Let $p$ represent the phase of $f$, and let $x\in[0,4)^2$ satisfy $x\equiv p+(1,1)\pmod4$. Then $x$, $x+\mathbf e_1$ and $x+\mathbf e_2$ lie in $[0,5)^2$. By \cref{prop:scaffold}(iii),(iv), the column of $x$ has residue $1$ relative to $p_1$ and column bit $0$, and the column of $x+\mathbf e_1$ has residue $2$ and column bit $1$. In the same way $\chi_2(f(x))=0$ and $\chi_2(f(x+\mathbf e_2))=1$. So $h=g|_{\operatorname{Anc}_0(x)}$ is a good $0$-anchor at $x$, and $h\in H$.

\emph{From depth $n$ to $n+1$.} Let $h_n\in H$ be a good $n$-anchor at $u$, with domain $\operatorname{Anc}_n(u)$. Let $U=[-L_3,31+L_3]^2$ and
\[
\Sigma=\operatorname{Anc}_n(u)\cup\bigcup_{x\in U}\operatorname{Port}_n(u+4^nx).
\]
Then $|\Sigma|\le3+\ell(2L_3+32)^2\le B\le K$, so $h_n$ extends to a member $g\in H$ with domain $\Sigma$. By \cref{scf:readout}, for every $x\in U$ the map $g$ reads a unique letter $\mathfrak c(x)$ on $\operatorname{Port}_n(u+4^nx)$. For $x,x+e\in U$, the edge certificate $\Gamma^{\mathrm{ed}}_n(u,x,e)$ with \cref{lem:transfer}, applied to the restriction of $g$ to its boundary, shows that $\mathfrak c(x)$ and $\mathfrak c(x+e)$ match in direction $e$. So $\mathfrak c$ is a legal pattern on $U$, and by \cref{scf:chart}(a) its restriction to $[0,32)^2$ is the restriction of a tiling $f$.

Apply \cref{prop:scaffold}(iii) to $f$ and then to its parent: there are $o\in[0,16)^2$ and a tiling $y'$ with $f(o+16m+i)=\sigma^2(y'(m))(i)$ for $i\in[0,16)^2$. By \cref{prop:scaffold}(iv) applied twice, every letter of $\sigma^2(\mathfrak b)$ in column $4$ has column bit $\chi_1(\sigma(\mathfrak b)(1,\cdot))=0$, and every letter in column $8$ has column bit $\chi_1(\sigma(\mathfrak b)(2,\cdot))=1$; rows $4$ and $8$ have row bits $0$ and $1$. Put $w=o+(4,4)$. The sites $w$, $w+4\mathbf e_1=o+(8,4)$ and $w+4\mathbf e_2=o+(4,8)$ lie in $o+[0,16)^2\subseteq[0,32)^2$. So
\[
\chi_1(\mathfrak c(w))=0\ne1=\chi_1(\mathfrak c(w+4\mathbf e_1)),\qquad \chi_2(\mathfrak c(w))=0\ne1=\chi_2(\mathfrak c(w+4\mathbf e_2)).
\]
By \cref{scf:ports}(c), $g(u+4^nx)=\sigma^n(\mathfrak c(x))(0,0)$ has the bits of $\mathfrak c(x)$. Put $u'=u+4^nw$. Then $\operatorname{Anc}_{n+1}(u')=\{u+4^nw,\ u+4^n(w+4\mathbf e_1),\ u+4^n(w+4\mathbf e_2)\}\subseteq\Sigma$, and $h_{n+1}=g|_{\operatorname{Anc}_{n+1}(u')}$ is a good $(n+1)$-anchor at $u'$. It lies in $H$ by restriction.
\end{proof}

The old anchor is dropped as soon as the new one is found, so no pins accumulate over the depths: each step uses one member of $H$ on at most $B$ sites. The anchor of depth $n$ may sit anywhere, and different families may place it differently. Together, \cref{prop:anchors} and \cref{scf:readout} give, in every $E_K$ family with $K\ge B$ and at every depth $n$, a member on $\operatorname{Anc}_n(u)\cup\operatorname{Port}_n(u)$, at most $\ell+3$ sites, that reads one root letter $\mathfrak a$; and every member that agrees with it there and is defined at a site of the block $u+[0,4^n)^2$ takes the value prescribed by the word $\sigma^n(\mathfrak a)$.
